\chapter{Cronograma e recursos}
\label{sec:cronograma}

\section{Etapas}

\begin{enumerate}
\item Protocolar os pedidos de acesso à informação ao Ministério da
Agricultura, à ADAGRI e à SEMACE.
\item Fixar as listas de municípios e assinar a pré-especificação.
\item Montar o painel de nascimentos e de óbitos fetais dos dois estados, de
2010 a 2024, e desde 2005 se os arquivos permitirem.
\item Calcular, só com o pré-período, o efeito mínimo detectável do desenho
principal e do desenho de calendário.
\item Medir a fração de domicílios próximos aos bananais.
\item Estimar: desfecho primário, diferença tripla, calendário, desfechos
secundários e mecanismos.
\item Segundo ensaio: preços, custo-benefício e discussão do instrumento.
\item Redação e defesa.
\end{enumerate}

\section{Cronograma}

O Quadro~\ref{qua:cronograma} distribui as etapas por trimestre. As datas são
uma proposta, a ajustar com o orientador e com o calendário do programa.

\begin{quadro}[ht]
\centering\small
\caption{Cronograma proposto}
\label{qua:cronograma}
\begin{tabular}{|l|c|c|c|c|c|}
\hline
Etapa & 2026/4 & 2027/1 & 2027/2 & 2027/3 & 2027/4 \\
\hline
1. Pedidos de acesso à informação & X & X & & & \\ \hline
2. Listas e pré-especificação     & X & X & & & \\ \hline
3. Painel de nascimentos e óbitos &   & X & & & \\ \hline
4. Cálculo de poder               &   & X & & & \\ \hline
5. Fração exposta                 &   & X & X & & \\ \hline
6. Estimação do primeiro ensaio   &   &   & X & X & \\ \hline
7. Segundo ensaio                 &   &   &   & X & X \\ \hline
8. Redação e defesa               &   &   & X & X & X \\ \hline
\end{tabular}
\fonte{Elaboração própria.}
\end{quadro}

\section{Recursos}

A pesquisa usa dados públicos e programas de código aberto (Python e R), e não
exige recursos além do tempo do pesquisador. Os pedidos de acesso à informação
são gratuitos. O código é mantido em repositório versionado.
